% concatenated from EtaBasis19112023.tex
%% 
%% Copyright 2007-2020 Elsevier Ltd
%% 
%% This file is part of the 'Elsarticle Bundle'.
%% ---------------------------------------------
%% 
%% It may be distributed under the conditions of the LaTeX Project Public
%% License, either version 1.2 of this license or (at your option) any
%% later version.  The latest version of this license is in
%%    http://www.latex-project.org/lppl.txt
%% and version 1.2 or later is part of all distributions of LaTeX
%% version 1999/12/01 or later.
%% 
%% The list of all files belonging to the 'Elsarticle Bundle' is
%% given in the file `manifest.txt'.
%% 

%% Template article for Elsevier's document class `elsarticle'
%% with numbered style bibliographic references
%% SP 2008/03/01
%%
%% 
%%
%% $Id: elsarticle-template-num.tex 190 2020-11-23 11:12:32Z rishi $
%%
%%
\documentclass[preprint,12pt]{elsarticle}

%% Use the option review to obtain double line spacing
%% \documentclass[authoryear,preprint,review,12pt]{elsarticle}

%% Use the options 1p,twocolumn; 3p; 3p,twocolumn; 5p; or 5p,twocolumn
%% for a journal layout:
%% \documentclass[final,1p,times]{elsarticle}
%% \documentclass[final,1p,times,twocolumn]{elsarticle}
%% \documentclass[final,3p,times]{elsarticle}
%% \documentclass[final,3p,times,twocolumn]{elsarticle}
%% \documentclass[final,5p,times]{elsarticle}
%% \documentclass[final,5p,times,twocolumn]{elsarticle}

%% For including figures, graphicx.sty has been loaded in
%% elsarticle.cls. If you prefer to use the old commands
%% please give \usepackage{epsfig}

%% The amssymb package provides various useful mathematical symbols
\usepackage{amssymb}
\usepackage{graphics}
%% The amsthm package provides extended theorem environments
\usepackage{amsmath}
\usepackage{amssymb}
\usepackage{amsthm}
\usepackage{xcolor}
\usepackage{enumerate}
\usepackage{subfig}
\usepackage{mathrsfs} 
\usepackage{amsthm}

%% The lineno packages adds line numbers. Start line numbering with
%% \begin{linenumbers}, end it with \end{linenumbers}. Or switch it on
%% for the whole article with \linenumbers.
%% \usepackage{lineno}


\theoremstyle{definition}
\newtheorem{definition}{Definition}

\theoremstyle{theorem}
\newtheorem{theorem}{Theorem}


\theoremstyle{lemma}
\newtheorem{lemma}{Lemma}

\theoremstyle{example}
\newtheorem{example}{Example}


\theoremstyle{remark}
\newtheorem{remark}{Remark}


\theoremstyle{prop}
\newtheorem{prop}{Proposition}




%\journal{Journal of Mathematical Analysis and Applications}

\begin{document}

\begin{frontmatter}

%% Title, authors and addresses

%% use the tnoteref command within \title for footnotes;
%% use the tnotetext command for theassociated footnote;
%% use the fnref command within \author or \address for footnotes;
%% use the fntext command for theassociated footnote;
%% use the corref command within \author for corresponding author footnotes;
%% use the cortext command for theassociated footnote;
%% use the ead command for the email address,
%% and the form \ead[url] for the home page:
%% \title{Title\tnoteref{label1}}
%% \tnotetext[label1]{}
%% \author{Name\corref{cor1}\fnref{label2}}
%% \ead{email address}
%% \ead[url]{home page}
%% \fntext[label2]{}
%% \cortext[cor1]{}
%% \affiliation{organization={},
%%             addressline={},
%%             city={},
%%             postcode={},
%%             state={},
%%             country={}}
%% \fntext[label3]{}

\title{A new Representation of $\alpha$-Bernstein Operators }

%% use optional labels to link authors explicitly to addresses:
%% \author[label1,label2]{}
%% \affiliation[label1]{organization={},
%%             addressline={},
%%             city={},
%%             postcode={},
%%             state={},
%%             country={}}
%%
%% \affiliation[label2]{organization={},
%%             addressline={},
%%             city={},
%%             postcode={},
%%             state={},
%%             country={}}

\author{Jamshid  Saeidian \corref{cor1}}
\ead{j.saeidian@khu.ac.ir}
\author{ Bahareh Nouri} 


%  
 \address{Faculty of Mathematical Sciences and Computer, Kharazmi University, No. 50,  Taleghani ave., Tehran, Iran.}
%% \ead[url]{home page}

 %\address{Faculty of Mathematical Sciences and Computer, Kharazmi University, No. 50,  Taleghani ave., Tehran, Iran.}
%% \fntext[label2]{}
\cortext[cor1]{Corresponding author}
%\author{Jamshid Saeidian \corref{cor1} and  Bahareh Nouri}
%
%\ead{j.saeidian@khu.ac.ir}
%\affiliation{organization={Faculty of Mathematical Sciences and Computer},%Department and Organization
%            addressline={Kharazmi University}, 
%            city={Tehran},
%         %   postcode={}, 
%         %   state={},
%            country={Iran}}

\begin{abstract}
%% Text of abstract
%The $\alpha$-Bernstein operators was first introduced in \textit{ [Chen, X., Tan, J., Liu, Z.,  Xie, J. (2017). Approximation of functions by a new family of generalized Bernstein operators. Journal of Mathematical Analysis and Applications, 450(1), 244-261.]}, and since then has been an inspiring force behind noumerous research endeavors. In this study, we present a new representation for this operators which is based on a recursive relation. This new representation gives us the opportunity to extract new properties of $\alpha$-Bernstein operators  and give alternative and simpler proofs for some theorems.

The $\alpha$-Bernstein operators were initially introduced in the paper by Chen, X., Tan, J., Liu, Z., Xie, J. (2017) titled "Approximation of Functions by a New Family of Generalized Bernstein Operators" (Journal of Mathematical Analysis and Applications, 450(1), 244-261). Since their introduction, these operators have served as a source of inspiration for numerous research endeavors. In this study,  we propose a novel  technique, founded on a recursive relation,  for constructing Bernstein-like bases. A special case of this new representation yields a novel portrayal of Chen's operators. This innovative representation enables the discovery of additional properties of $\alpha$-Bernstein operators and facilitates alternative and more straightforward proofs for certain theorems.
\end{abstract}

%%Graphical abstract
%\begin{graphicalabstract}
%\includegraphics{grabs}
%\end{graphicalabstract}

%%Research highlights
%\begin{highlights}
%\item Research highlight 1
%\item Research highlight 2
%\end{highlights}

\begin{keyword}
%% keywords here, in the form: keyword \sep keyword
%% PACS codes here, in the form: \PACS code \sep code
%% MSC codes here, in the form: \MSC code \sep code
%% or \MSC[2008] code \sep code (2000 is the default)
Bernstein-like operator \sep  Generalized operators \sep Positive linear operator \sep Shape-preserving approximation \sep Voronovskaja-type theorem

\MSC[2020]  41A25 \sep 26A15 \sep  47A58 \sep 41A10

\end{keyword}

\end{frontmatter}

%% \linenumbers

%% main text
\section{Introduction}  \label{}

%
%The Bernstein polynomials, rooted in the fertile ground of approximation theory, stand as an enduring pillar of mathematical research and practice. These polynomials, introduced by S.N. Bernstein in 1912, not only define a fundamental class of algebraic polynomials but also wield a profound influence in mathematical analysis. It was through the lens of Bernstein polynomials that S.N. Bernstein provided the first constructive proof of Weierstrass' approximation theorem, a seminal contribution that continues to reverberate through the annals of mathematical history \cite{farouki2012bernstein}.
%
%Defined as:
%
%\[p_{n,k}(z) ={n \choose k} z^k(1-z)^{n-k}, ~~k=0,...,n ~~ \& ~~ n\in \mathbb{N},\]
%
%these polynomials are instrumental in approximating functions over the interval \([0,1]\). The core of this approximation is encapsulated in the formula:
%
%\[ \mathcal{B}_n(f;z) =\sum_{k=0}^n p_{n,k}(z)f(\frac{k}{n}),  ~~~ z\in [0,1],\]
%
%where a weighted sum of function values at equidistant points on the interval offers a versatile tool for approximation.
%
%The study of Bernstein polynomials and operators transcends mere mathematical curiosity; it has led to a vibrant and dynamic field of research. A testament to their significance, Lorentz's book stands as a beacon in the study of Bernstein polynomials, illuminating their properties and applications \cite{lorentz2012bernstein}. Yet, the influence of Bernstein polynomials extends far beyond their intrinsic qualities. Their form has served as a source of inspiration for mathematicians, spurring the development of an extensive array of approximation operators.
%
%Since their inception, Bernstein polynomials have been the muse for a diverse assortment of approximation methods, each designed to cater to distinct challenges and problem domains. The family of adaptations includes Schurer polynomials, Kantorovich polynomials, Stancu polynomials, $q$-Bernstein polynomials, Durrmeyer polynomials, Favard-Sz\'asz-Mirakjan operators, Baskakov operators, and many others \cite{gal2009approximation, gupta2013approximation, govil2013certain}. Recent strides in this area can be traced through references such as \cite{ozger2020approximation, kajla2021blending, mohiuddine2023order}.
%
%The development and analysis of Bernstein-type operators have been guided by two primary objectives: preserving the forms of diverse functions, including polynomials and exponentials, and upholding their shape-preserving properties \cite{cardenas2006shape, aral2018bernstein}.
%
%%In 2003, King \cite{king2003positive} introduced a sequence of Bernstein-type operators notable for their capability to preserve the function \(x^2\). This pioneering study laid the foundation for subsequent research efforts, which further refined and extended these developments \cite{gonska2009general,cardenas2011bernstein}.
%
%Recent years have witnessed a resurgence of interest in the development of Bernstein-type operators. In 2017, Chen et al. \cite{chen2017approximation} proposed an innovative modification of Bernstein operators, centered around the concept of $\alpha$-Bernstein polynomials. This groundbreaking work has ignited substantial interest and inspired further exploration into diverse families of Bernstein-type operators \cite{mohiuddine2017construction, kajla2018blending}.





The Bernstein polynomials, rooted in the fertile ground of approximation theory, stand as an enduring pillar of mathematical research and practice. These polynomials, introduced by S.N. Bernstein in 1912, not only define a fundamental class of algebraic polynomials but also wield a profound influence in mathematical analysis. It was through the lens of Bernstein polynomials that S.N. Bernstein provided the first constructive proof of Weierstrass' approximation theorem, a seminal contribution that continues to reverberate through the annals of mathematical history \cite{farouki2012bernstein}.

Defined as:
\[B_{n,k}(z) ={n \choose k} z^k(1-z)^{n-k}, ~~k=0,...,n ~~ \& ~~ n\in \mathbb{N},\]
these polynomials are instrumental in approximating functions over the interval \([0,1]\). The core of this approximation is encapsulated in  Bernstein operators with the formula:
\[ \mathcal{B}_n(f;z) =\sum_{k=0}^n B_{n,k}(z)f(\frac{k}{n}),  ~~~ z\in [0,1],\]
where a weighted sum of function values at equidistant points on the interval offers a versatile tool for approximation.

In 2017, Chen et al. \cite{chen2017approximation} proposed an innovative modification of Bernstein operators, centered around the concept of $\alpha$-Bernstein polynomials as follow:
\begin{definition}\cite{chen2017approximation}
The $ \alpha $-Bernstein polynomials  of degree $n$, $B_{n,k}^{(\alpha)}(z)$, are defined by $B_{1,0}^{(\alpha)}(z)=1-z$, $B_{1,1}^{(\alpha)}(z)=z$ and
\begin{equation} 
B_{n,k}^{(\alpha)}(z) = \left[ \binom{n-2}{k}(1-\alpha)z + \binom{n-2}{k-2}(1-\alpha)\left(1-z\right) + \binom{n}{k}\alpha  z(1-z) \right]z^{k-1}(1-z)^{n-k-1}   \label{chenformula}
\end{equation}
where $n \geq 2$, $z \in[0, 1]$ and $\binom{m}{l}$ denotes  the binomial coefficients.
\end{definition}


Then, the corresponding $ \alpha $-Bernstein operators are defined to read
\[ \mathcal{B}_n^{\alpha}(f;z) =\sum_{k=0}^n B_{n,k}^{\alpha}(z)f(\frac{k}{n}),  ~~~ z\in [0,1],\] 
for any $f \in  C[0,1]$, i.e. continuous functions.

In recent years, the field of approximation theory has experienced noteworthy advancements, propelled in part by the pioneering work of Chen and his colleagues. Their contributions to the theory of Bernstein-like operators have not only broadened our comprehension but also acted as a catalyst for subsequent exploration and refinement. This seminal work has sparked significant interest and inspired further investigations into various families of Bernstein-like operators \cite{mohiuddine2017construction, kajla2018blending, mohiuddine2020approximation, ozger2020approximation}.


%In recent years, the field of approximation theory  has witnessed significant advancements, spurred in part by this groundbreaking work of Chen and his colleagues. Their contributions to the theory of Bernstein-like operators have not only expanded our understanding but have also become a catalyst for further exploration and refinement. This seminal  work has ignited substantial interest and inspired further exploration into diverse families of Bernstein-type operators \cite{mohiuddine2017construction, kajla2018blending, mohiuddine2020approximation, ozger2020approximation}.


The $\alpha$-Bernstein operators are commonly identified as generalized Bernstein operators in the literature, and researchers have further developed the underlying concept of this modification to generalize other prominent operators in the field.
Among these extensions, we can mention operators linked to Schurer polynomials, Kantorovich polynomials, Stancu polynomials, $q$-Bernstein polynomials, Durrmeyer polynomials, as well as other operators like Favard-Szász-Mirakjan and Baskakov, these  extensions, as well as other related works could be traced in \cite{    mohiuddine2017construction, ozger2020approximation,  mohiuddine2021approximation, aral2019parametric, cheng2023construction, kajla2019generalized}.


Inspired by Chen's research, this paper presents a novel framework for generating Bernstein-like basis functions. This innovative structure can be viewed as an extension of Bernstein polynomials, aligning with these classical polynomials and their $\alpha$-Bernstein modifications  \cite{chen2017approximation} in specific scenarios. This framework enables us to  introduce a recursive  representation of the  $\alpha$-Bernstein operators, a representation that promises to enhance our comprehension of its underlying principles and applications. By revisiting the operator's structure through a fresh lens, we aim to offer a unique perspective on its mathematical properties and computational implications.

%Motivated by the desire to deepen our understanding of Chen's operator, we unveil a new parametrization that not only refines the operator's representation but also facilitates alternative approaches to the proof of key theorems. Morover, new properties are achieved 

Our primary focus lies in presenting this new representation and demonstrating its utility through  alternative  proofs for selected theorems associated with $\alpha$-Bernstein operators.
Moreover, it provides us with the opportunity to uncover some additional properties of $\alpha$-Bernstein operators.
% This approach not only contributes to the ongoing development of Bernstein-like operators but also establishes a foundation for future investigations in approximation theory and CAGD.


%The remainder of this paper unfolds as follows: Section 2 outlines the new representation and verifies the equvalency of the new presentation with Chen's formula. Section 3 studies  the  properties of the new presentation in light of the Bernstein-type basis.  In Section 4, in one hand two new properties of the   $\alpha$-Bernstein operators  are extracted and in the sequel  we present alternative proofs for key theorems, showcasing the practical implications of our proposed representation. Finally, Section 5 offers concluding remarks and suggests avenues for future research.

The structure of the rest of this paper is as follows: Section 2 delineates the novel representation and validates its equivalence with Chen's formula. Section 3 examines the properties of this new presentation within the context of the Bernstein-like bases. In Section 4, on one hand, two novel properties of the $\alpha$-Bernstein operators are derived, and subsequently, we provide alternative proofs for key theorems, illustrating the practical implications of our proposed representation. Finally, Section 5 presents concluding remarks and proposes potential avenues for future research.


\section{A general structure for constructing Bernstein-like bases}

The idea of recursively constructing Bernstein-like basis functions has already been employed by researchers,
 Yan and Liang \cite{yan2011extension} utilized this approach to construct a Bezier-like model and  Li \cite{li2018novel} extended the idea to define a parameter-based Bernstein-like basis, suitable for constructing Bezier-type curves. Additionally, Bibi et al. \cite{bibi2020geometric} introduced the so-called \textit{generalized hybrid trigonometric Bernstein basis}, and Ameer and his colleagues \cite{ameer2022novel} developed generalized Bernstein-like basis functions. Both of these studies further contribute to the recursive construction of Bernstein-like bases.
%Yan and Liang \cite{yan2011extension} used the  idea for constructing   a Bezier-like  model.
% Li \cite{li2018novel} utilized  the idea to define a parameter based Bernstein-type basis which could be used to construct Bezier-type curves. The so called \textit{generalized hybrid trigonometric Bernstein basis} constructed by Bibi et al. \cite{bibi2020geometric}  and 
%the generalized Bernstein-like basis functions founded by Ameer and his colleagues  \cite{ameer2022novel} aretwo more studies which are based on recursive construction of Bernstein-like bases.

The idea is straightforward: identify an appropriate set of initial bases, such as $\{F_{2,0},F_{2,1},F_{2,2}\}$, and generate higher-order bases through the established recursive relation of Bernstein polynomials \cite{farin2002curves}:
%The idea is so simple, find a suitable set of  initial bases, say for example $\{F_{2,0},F_{2,1},F_{2,2}\}$  and construct  higher order bases using the well-known recursive relation of Bernstein polynomials
  \begin{equation}
F_{n,i}(z) = (1-z)F_{n-1,i}(z) + zF_{n-1,i-1}(z) \quad \text{for } z \in [0,1], \quad i = 0, 1, 2, \ldots, n. \label{rec-rel}
\end{equation}
%The number of initial bases, which we prefer to call them \textit{starting basis} from then on, could be any natural number, Yan and Liang \cite{yan2011extension} and Ameer et al.  \cite{ameer2022novel}, both, used  three polynomials of degree 4,  Li \cite{li2018novel} sets  four polynomials of degree 3 as starting basis while Bibi et al \cite{bibi2020geometric} employed three hybrid trigonometric functions.
%It is required from the  starting basis functions to share the  fundamental properties of Bernstein bases  as many as possible.

The count of initial bases, hereafter referred to as the \textit{starting basis}, can be any natural number. Yan and Liang \cite{yan2011extension} and Ameer et al. \cite{ameer2022novel} both employed three polynomials of degree 4, Li \cite{li2018novel} selected four polynomials of degree 3 as the starting basis, and Bibi et al. \cite{bibi2020geometric} utilized three hybrid trigonometric functions. It is essential for the starting basis functions to exhibit as many fundamental properties of Bernstein bases as possible.

We stick to this technique and propose  a general structure to define a family of Bernstein-like bases.

\begin{definition} \label{defphi}
The starting Bernstein-like basis functions of order two are  defined  for $z \in [0,1]$ as 
\begin{eqnarray} 
g_{2,0}(z)&=&az+b+\varphi(z),\nonumber \\
g_{2,1}(z)&=&1-2b-a-2\varphi(z), \label{formula-phi}\\
g_{2,2}(z)&=&-az+b+a+\varphi(z),\nonumber
\end{eqnarray} 
where $a,b \in \mathbb{R}$  and the real valued function $\varphi$ must satisfy some conditions. The higher order Bernstein-like bases are  generated by recursive relation (\ref{rec-rel}).
\end{definition}

\begin{lemma}\label{lem1} 
 The Bernstein-like basis functions of order \(n \geq 2\), as defined in Definition \ref{defphi}, satisfy the partition of unity property, expressed as \( \sum_{i=0}^{n}g_{n,i}(z)=1 \). Additionally, if the function \(\varphi\) is constrained to satisfy \(\varphi(1-z)=\varphi(z)\), then a symmetry feature emerges: \( g_{n,i}(z)=g_{n,n-i}(1-z) \) for \( i = 0, 1, 2, \ldots , n \).

\end{lemma}
%\begin{lemma}
%If the function $\varphi$ satisfies  $\varphi(1-z)=\varphi(z)$, then the Bernstein-like basis functions of order $n \geq 2$, defined in Definition \ref{defphi}, hold the following properties  
%\begin{itemize}
%\item [\rm{(a)}]  Partition of unity: $ \sum_{i=0}^{n}g_{n,i}(z)=1 $.
%\item [\rm{(b)}] Symmetry: $ g_{n,i}(z)=g_{n,n-i}(1-z) $ for $ i = 0, 1, 2, \ldots , n $. 
%\end{itemize}
%\end{lemma}

One can impose further restrictions on the  function $\varphi$ and the real values $a$ and $b$ in order to satisfy more fundamental properties of classical Bernstein polynomials. In fact,  the non-negativity and end-point coincidence with Bernstein bases are two unavoidable features.

\begin{example}
When  we set $a=-1$, $b=0$ and $\varphi(z)=z^2-z+1$, the  polynomials in equation (\ref{formula-phi}) reduce to the   Bernstein polynomials of degree 2, thus Definition \ref{defphi} constructs the classical Bernstein polynomials.
\end{example}

\begin{example}
Authors in \cite{nouri2022class} presented a special case of (\ref{formula-phi}) where they used the  function $\varphi (z)=\sqrt{\left ( 1-\nu  \right )\left ( z^{2}-z \right )+\frac{1}{4}}$, along with $a=-1$ and $b=\frac{1}{2}$. 
%It is verified that the so-called $sq$-basis functions satisfy the non-negativity and end-point interpolation properties as well. Moreover, the corresponding operators share the most common features of the classical Bernstein operators \cite{nouri2023family}.
It is verified that the so-called $sq$-basis functions also satisfy the non-negativity and end-point interpolation properties. Furthermore, the corresponding operators exhibit the most common features of the classical Bernstein operators \cite{nouri2023family}.
\end{example}





\section{New representation of the $\alpha$-Bernstein polynomials}

Setting $a=-1$, $b=1$  and $\varphi(z)=\alpha (z^2-z)$, in equation (\ref{formula-phi}), leads us to a set of starting basis functions as follows:
\begin{eqnarray} 
F_{2,0}(z)&=&1-z+\alpha \left ( z^{2}-z \right ),\nonumber \\
F_{2,1}(z)&=&-2\alpha \left ( z^{2}-z \right ), \label{formula1}\\
F_{2,2}(z)&=&z+\alpha \left ( z^{2}-z \right ).\nonumber
\end{eqnarray} 
Now employing the recurrence relation (\ref{rec-rel}), one constructs a set of basis functions which its equivalence with $\alpha$-Bernstein polynomials of Chen \cite{chen2017approximation} are easily verified, thus we have:
\begin{theorem}

The procedure presented in Definition \ref{defphi}, along with starting basis functions (\ref{formula1}), generates the $\alpha$-Bernstein polynomials.
\end{theorem}



The recursive definition of the $\alpha$-Bernstein polynomials provides an opportunity to derive practical computational formulas, one of which is articulated in the following lemma.

\begin{lemma}\label{lemsigma}
The  $\alpha $-Bernstein polynomials, defined by equations (\ref{rec-rel}) and (\ref{formula1}), can be derived through the following relation:
\begin{eqnarray}\label{sigma-eq}
F_{n,i}(z) =\sum_{j=0}^{n-2}B_{n-2,j}(z)F_{2,i-j}(z),  ~~i=0,1,\cdots,n,
\end{eqnarray}
where  $ B_{n-2,j}  $  are the classical Bernstein polynomials of order $n-2$.
\end{lemma}
\begin{proof}
We utilize the method of induction to establish the validity of the result, initiating with the base case  $n = 3$:
\begin{eqnarray*}
F_{n,0}(z) &=&\sum_{j=0}^{1}B_{1,j}(z)F_{2,-j}(z)=B_{1,0}(z)F_{2,0}(z)=(1-z)F_{2,0}(z),\\
F_{n,1}(z) &=&\sum_{j=0}^{1}B_{1,j}(z)F_{2,1-j}(z)=B_{1,0}(z)F_{2,1}(z)+B_{1,1}(z)F_{2,0}(z)=(1-z)F_{2,1}(z)+zF_{2,0}(z),\\
F_{n,2}(z) &=&\sum_{j=0}^{1}B_{1,j}(z)F_{2,2-j}(z)=B_{1,0}(z)F_{2,2}(z)+B_{1,1}(z)F_{2,1}(z)=(1-z)F_{2,2}(z)+zF_{2,1}(z),\\
F_{n,3}(z) &=&\sum_{j=0}^{1}B_{1,j}(z)F_{2,3-j}(z)=B_{1,1}(z)F_{2,2}(z)=zF_{2,2}(z),\\
\end{eqnarray*}
the desired outcome is achieved.


Now, assuming the validity of the relation for $ n-1 $, we obtain:
\begin{eqnarray*}
F_{n-1,i}(z) =\sum_{j=0}^{n-3}B_{n-3,j}(z)F_{2,i-j}(z), ~~ i=0,1,\cdots,n-1,
\end{eqnarray*}
In the next step, we aim to establish the validity of the relation for  $n$. Utilizing the recursive relation (\ref{rec-rel}), for  $i=0,1,\cdots,n$, we have:
\begin{eqnarray*}
F_{n,i}(z)& =& (1-z)F_{n-1,i}(z) + zF_{n-1,i-1}(z)\\
&=&(1-z)\sum_{j=0}^{n-3}B_{n-3,j}(z)F_{2,i-j}(z)+ z\sum_{j=0}^{n-3}B_{n-3,j}(z)F_{2,i-1-j}(z)\\
&=&(1-z)B_{n-3,0}(z)F_{2,i-0}(z)+\sum_{j=1}^{n-3}(1-z) B_{n-3,j}(z)F_{2,i-j}(z)\\
&&~~~~~~~+ \sum_{j=1}^{n-3}zB_{n-3,j-1}(z)F_{2,i-j}(z)+zB_{n-3,n-3}(z)F_{2,i-(n-2)}(z)\\
&=&(1-z)B_{n-3,0}(z)F_{2,i-0}(z)+\sum_{j=1}^{n-3}[(1-z) B_{n-3,j}(z)+zB_{n-3,j-1}(z) ]F_{2,i-j}(z)\\
  &&    ~~~~~~~+zB_{n-3,n-3}(z)F_{2,i-(n-2)}(z)\\
&=&B_{n-2,0}(z)F_{2,i-0}(z)+\sum_{j=1}^{n-3}B_{n-2,j}(z)F_{2,i-j}(z)+B_{n-2,n-2}(z)F_{2,i-(n-2)}(z)\\
&=&\sum_{j=0}^{n-2}B_{n-2,j}(z)F_{2,i-j}(z).
\end{eqnarray*}


\end{proof}


\begin{remark}
For any given value of \(n\), the right-hand side of (\ref{sigma-eq}) has at most  three non-zero terms, providing a computational advantage. Moreover, relation (\ref{sigma-eq}) can be represented in a matrix form as follows:
\begin{eqnarray*}
 \begin{bmatrix}
F_{n,0}(z)\\ 
F_{n,1}(z)\\ 
\vdots \\
F_{n,n}(z)
\end{bmatrix}=
\begin{bmatrix}
B_{n-2,0}(z) & 0 &0 \\ 
B_{n-2,1}(z) & B_{n-2,0}(z) & 0 \\ 
B_{n-2,2}(z) & B_{n-2,1}(z) & B_{n-2,0}(z) \\ 
B_{n-2,3}(z) & B_{n-2,2}(z) & B_{n-2,1}(z)  \\ 
\vdots  &  \vdots & \vdots \\ 
B_{n-2,n-3}(z) & B_{n-2,n-4}(z)  & B_{n-2,n-5}(z) \\ 
B_{n-2,n-2}(z) & B_{n-2,n-3}(z)  & B_{n-2,n-3}(z) \\ 
0 & B_{n-2,n-2}(z) & B_{n-2,n-3}(z) \\ 
 0& 0 & B_{n-2,n-2}(z)
\end{bmatrix}%_{(n+1)\times 3}
\begin{bmatrix}
F_{2,0}(z)\\ 
F_{2,1}(z)\\ 
F_{2,2}(z)
\end{bmatrix}.%_{3\times 1}
\end{eqnarray*}
%
%\begin{eqnarray*}
%F_{n,i}(z) = \begin{bmatrix}
%B_{n-2,0}(z),
%B_{n-2,1}(z),
%B_{n-2,2}(z) ,
%\cdots,
%B_{n-2,n-2}(z)
%\end{bmatrix}
%\underbrace{
%\begin{bmatrix}
%1 & 0 &0 \\ 
%1 & 1& 0 \\ 
%1 &1 & 1 \\ 
%1 & 1& 1  \\ 
%\vdots  &  \vdots & \vdots \\ 
%1& 1  & 1\\ 
%1 &1  & 1 \\ 
%0 &1& 1\\ 
% 0& 0 & 1
%\end{bmatrix}}_{T_{(n-1)\times 3}}
%\begin{bmatrix}
%F_{2,0}(z)\\ 
%F_{2,1}(z)\\ 
%F_{2,2}(z)
%\end{bmatrix}.%_{3\times 1}
%\end{eqnarray*}
\end{remark}


From the proof of Lemma \ref{lemsigma}, the following fact is straightforward:

\begin{lemma} \label{rec-lem}
For any set of basis functions $\{F_{n,i}\}_{i=0}^n$, which are constructed using the recursive relation (\ref{rec-rel}), any basis of order $n$ could be represented in terms of the basis elements of order $m <n$, i.e.
\begin{eqnarray} \label{sigma}
F_{n,i}(z) =\sum_{j=0}^{n-m}B_{n-m,j}(z)F_{m,i-j}(z), ~~~ i=0,\cdots, n.
\end{eqnarray}

\end{lemma}

\begin{remark}
%Formula (\ref{sigma}), reduces the computational complexity in computing the basis terms. 
 Lemma \ref{rec-lem} provides a universal formula for expressing the basis functions in terms of  the starting basis functions, irrespective of the number of starting bases employed.
\end{remark}









The forthcoming result demonstrates that the $\alpha$-Bernstein polynomials possess many of the properties of classical Bernstein polynomials, making them well-suited for function approximation.

\begin{theorem}\label{thm1} 
The $\alpha$-Bernstein polynomials  derived from equations  (\ref{rec-rel}) and (\ref{formula1}) exhibit the following  properties:
\begin{itemize}
\item [\rm{(a)}] Nonnegativity: $ F_{n,i}(z) \geq 0$  for $ i = 0, 1, 2,\ldots , n$.
\item [\rm{(b)}] Partition of unity: $ \sum_{i=0}^{n}F_{n,i}(z)=1 $.
\item [\rm{(c)}] Symmetry: $ F_{n,i}(z)=F_{n,n-i}(1-z) $ for $ i = 0, 1, 2, \ldots , n $. 
\item [\rm{(d)}] End-point values:
\begin{eqnarray}
F_{n,i}(0)= 
\begin{cases} 
1, & i=0,
\\ 
0, & i\neq 0,
\end{cases}
\hspace{0.7cm}
F_{n,i}(1)= 
\begin{cases} 
1, & i=n,
\\ 
0, & i\neq n.
\end{cases}
\label{end-val}
\end{eqnarray}
%\item[(e)]   \textcolor{red}{Unimodality?}
\end{itemize}
\end{theorem}






\begin{proof}
%The process of establishing the proof of this theorem relies on the utilization of mathematical induction.
%Mathematical induction is used to prove this theorem.
The proof of this theorem relies on the application of the  mathematical induction.
\begin{itemize}
\item [\rm{(a)}]
Non-negativity is obvious by  (\ref{rec-rel}) and  (\ref{formula1}).
\item [\rm{(b)}]
It is  already stated in Lemma \ref{lem1}.
%In Equation (\ref{formula1}), we derive $ \sum_{i=0}^{2}F_{2,i}(z)=1 $. Now, let's assume that the equality holds for $k$. Using the recurrence relation (\ref{rec-rel}) and the inductive hypothesis, we can then deduce 
%From (\ref{1}),we obtain $ \sum_{i=0}^{2}F_{2,i}(z)=1 $. 
%Now suppose that the equality is true for $k$, then
%according to recurrence relation(\ref{2}) and the inductive hypothesis, we
%obtain 
%$$
%\sum_{i=0}^{k+1}F_{k+1,i}(t) = (1 - z)\sum_{i=0}^{k+1}F_{k,i}(t) + z\sum_{i=0}^{k+1}F_{k,i-1}(t) =1-t+t=1.
%$$
\item [\rm{(c)}]
%The $\alpha$-starting basis functions are symmetrical. Assume that the $\alpha$-basis functions of order $k$ are symmetrical, 
%then from this inductive hypothesis and
%recurrence relation (\ref{2}), one has 
The starting basis functions exhibit symmetry. Let's assume that the $\alpha$-Bernstein polynomials of order $k$  hold  symmetry. Now, by considering this inductive hypothesis along with the recurrence relation (\ref{rec-rel}), we can conclude that
\begin{eqnarray*}
F_{k+1,i}(1-z)&=&(1-(1-z))F_{k,i}(1-z)+(1-z)F_{k,i-1}(1-z)\\
&=&zF_{k,k-i}(z)+(1-z)F_{k,k-i+1}(z)=F_{k+1,k-i+1}(z).
\end{eqnarray*}
\item [\rm{(d)}]
%By a simple deduction from (\ref{1}) , it conclude that the results in (\ref{rr1}) hold for $n=2$.
Based on a straightforward deduction from (\ref{formula1}), it can be concluded that the results in (\ref{end-val}) apply to the case when $n=2$.
\\
%Suppose that the properties at the endpoints hold for the $\alpha$-basis functions of order $ k$, so we have
Let's assume that the properties at the endpoints are satisfied by the $\alpha$-Bernstein polynomials of order $k$. Consequently, by utilizing the inductive hypothesis and equation (\ref{rec-rel}), we can conclude that:
\begin{eqnarray}
F_{k+1,i}(0)&=&F_{k,i}(0)= 
\begin{cases} 
1 & i=0,
\\ 
0, & i\neq 0,
\end{cases} \nonumber
\\ %\hspace{0.01cm}
F_{k+1,i}(1)&=&F_{k,i-1}(1)= 
\begin{cases} 
1 & i=k+1,
\\ 
0, & i\neq k+1.
\end{cases} \nonumber
\end{eqnarray}
%which can be obtained by the inductive hypothesis and (\ref{2}). 


\end{itemize}
\end{proof}
In Figure \ref{fig1}, an illustration of the $\alpha$-Bernstein polynomials  is presented. These functions are generated for varying values of $ \alpha $ (including $ 0,0.2, 0.4,0.6,0.8$ and $ 1 $), while the values of $ n $ are set to  $2, 3,$ and $4$.

%Figure \ref{fig1} shows the $ \alpha $-basis functions generated for  $ n=2,3,4$ and value of $ \alpha$ are set to $\beta = 0,0.2, 0.4,0.6,0.8,1$.
\begin{figure}[ht]
\centering
\subfloat 
{\includegraphics[width=6.5cm]{fig/pic1}}
\\
\subfloat
{\includegraphics[width=6.5cm]{fig/pic2}}
\hspace*{0.1mm}
\subfloat
{\includegraphics[width=6.5cm]{fig/pic3}}
\caption{The $\alpha$-Bernstein polynomials  by adjusting $ \alpha $ from  $ 0, 0.2,\ldots,1$.} 
\label{fig1}
\end{figure}




\section{New results and alternative proofs}

%This section is devoted to demonstrate the advantages of the new representation of  $\alpha$-Bernstein operators,  we present  a new relation for computing the moments and illustrate  a Vornovskaya type theorem. Afterwards the shape-preserving properties of this operators are verified in a completely different strategies.

This section is dedicated to showcasing the benefits of the new representation of $\alpha$-Bernstein operators. We introduce a new formula for computing the moments and present an analogue of Vornovskaja's theorem. Subsequently, we verify the shape-preserving properties of these operators using entirely different strategies.



\begin{theorem}
For all $j \in \mathbb{N}\cup \left \{ 0 \right \}, n \in \mathbb{N}$, $\alpha \in [0,1]$ and $  z \in \left [ 0,1 \right ] $, we have the recurrence
formula:
\begin{eqnarray*}
\mathcal{B}_n^{\alpha}\left ( z^{j+1};z \right )=\left ( 1-\frac{1}{n} \right )^{j+1}\left [ \left ( 1-z \right )\mathcal{B}_{n-1}^{\alpha}\left ( z^{j+1};z \right )+z\sum_{k=0}^{j+1} \binom{j+1}{k}\dfrac{1}{\left ( n-1 \right )^{k}} \mathcal{B}_{n-1}^{\alpha}\left ( z^{j+1-k};z \right )\right ].
\end{eqnarray*}
\end{theorem}
\begin{proof}
Employing equation (\ref{rec-rel})  we derive:
\begin{eqnarray*}
\mathcal{B}_n^{\alpha}\left ( z^{j+1};z \right )&=& \sum_{i=0}^{n}\left ( \frac{i}{n} \right )^{j+1}F_{n,i}(z)\\
&=&\sum_{i=0}^{n}\left ( \frac{i}{n} \right )^{j+1}\left [ \left ( 1-z \right )F_{n-1,i}(z)+zF_{n-1,i-1}(z) \right ]\\
&=&\frac{\left ( n-1 \right )^{j+1}}{n^{j+1}}\left [ \left ( 1-z \right )\sum_{i=0}^{n}\left ( \frac{i}{n-1} \right )^{j+1}F_{n-1,i}(z)+z\sum_{i=0}^{n}\left ( \frac{i}{n-1} \right )^{j+1}F_{n-1,i-1}(z) \right ]\\
&=&\left ( 1-\frac{1}{n} \right )^{j+1}\left [ \left ( 1-z \right )\mathcal{B}_{n-1}^{\alpha}\left ( z^{j+1};z \right )+z\sum_{i=0}^{n}\left ( \frac{i-1+1}{n-1} \right )^{j+1}F_{n-1,i-1}(z)  \right ]\\
&=&\left ( 1-\frac{1}{n} \right )^{j+1}\left [ \left ( 1-z \right )\mathcal{B}_{n-1}^{\alpha}\left ( z^{j+1};z \right )+z\sum_{k=0}^{j+1} \binom{j+1}{k}\dfrac{1}{\left ( n-1 \right )^{k}} \sum_{i=0}^{n-1}\left ( \frac{i}{n-1} \right )^{j+1-k}F_{n-1,i}(z) \right ]\\
&=&\left ( 1-\frac{1}{n} \right )^{j+1}\left [ \left ( 1-z \right )\mathcal{B}_{n-1}^{\alpha}\left ( z^{j+1};z \right )+z\sum_{k=0}^{j+1} \binom{j+1}{k}\dfrac{1}{\left ( n-1 \right )^{k}} \mathcal{B}_{n-1}^{\alpha}\left ( z^{j+1-k};z \right )\right ].
\end{eqnarray*}
We utilized the outcome of the following relation in the above calculations.
\begin{eqnarray*}
\sum_{i=0}^{n}\left( \frac{i-1+1}{n-1} \right)^{j+1}F_{n-1,i-1}(z) 
&=&\sum_{i=0}^{n}\dfrac{1}{\left ( n-1 \right )^{j+1}}\left [\sum_{k=0}^{j+1}\binom{j+1}{k}\left ( i-1 \right )^{j+1-k}  \right ]F_{n-1,i-1}(z) \\
&=&\sum_{i=0}^{n} \left[\sum_{k=0}^{j+1}\binom{j+1}{k}\dfrac{1}{\left( n-1 \right)^{k}}\left( \dfrac{i-1}{n-1} \right)^{j+1-k}  \right]F_{n-1,i-1}(z) \\
&=&\sum_{k=0}^{j+1} \binom{j+1}{k}\dfrac{1}{\left( n-1 \right)^{k}} \sum_{i=0}^{n-1}\left( \frac{i}{n-1} \right)^{j+1-k}F_{n-1,i}(z). 
\end{eqnarray*}
\end{proof}




Next, we  present a Grüss-Voronovskaja-type theorem for the $\alpha$-Bernstein operators, utilizing the methodology outlined in \cite{gal2014gr}.
\begin{theorem}
Let $f, h \in C^{2}[0, 1]$,  then for any $z \in [0, 1]$, we have 
\begin{equation}
\underset{n\rightarrow \infty}{lim}~~n \left [\mathcal{B}_n^{\alpha}(fh; z) -\mathcal{B}_n^{\alpha}(f; z)\mathcal{B}_n^{\alpha}(h; z)  \right ]= z(1-z){f}'(z){h}'(z),
\end{equation}
where $ \alpha\in [0,1]$.
\end{theorem}
\begin{proof}
It is straightforward to express
\begin{eqnarray*}
&&\mathcal{B}_n^{\alpha}(fh; z) -\mathcal{B}_n^{\alpha}(f; z)\mathcal{B}_n^{\alpha}(h; z)\\
&=&\mathcal{B}_n^{\alpha}(fh; z)-f\left ( z \right )h\left ( z \right )+f\left ( z \right )h\left ( z \right ) -\mathcal{B}_n^{\alpha}(f; z)\mathcal{B}_n^{\alpha}(h; z)+\mathcal{B}_n^{\alpha}(f; z)h\left ( z \right )-\mathcal{B}_n^{\alpha}(f; z)h\left ( z \right )\\
&=&\left [ \mathcal{B}_n^{\alpha}(fh; z)-f\left ( z \right )h\left ( z \right ) \right ]+\left [ h\left ( z \right ) \Big ( f(z)-\mathcal{B}_n^{\alpha}(f; z)\Big)  \right ]+\left [ \mathcal{B}_n^{\alpha}(f; z) \Big ( h(z)-\mathcal{B}_n^{\alpha}(h; z)\Big)  \right ].
\end{eqnarray*}
Considering corresponding Korovkin \cite{korovkin1953convergence} and Vornovskaja \cite{voronovskaja1932determination}  theorems from  the original paper \cite{chen2017approximation}, we derive the following
\begin{eqnarray*}
&&
\underset{n\rightarrow \infty}{lim} n \left[\mathcal{B}_n^{\alpha}(fh; z) -\mathcal{B}_n^{\alpha}(f; z)\mathcal{B}_n^{\alpha}(h; z)\right]\\
&=&
\underset{n\rightarrow \infty}{lim}n \left [ \mathcal{B}_n^{\alpha}(fh; z)-f\left ( z \right )h\left ( z \right ) \right ]+\left [ h\left ( z \right )
 \Big (\underset{n\rightarrow \infty}{lim} n [f(z)-\mathcal{B}_n^{\alpha}(f; z)]\Big)  \right ]\\
 &&~~~~+\left [ 
\underset{n\rightarrow \infty}{lim} \mathcal{B}_n^{\alpha}(f; z) 
 \Big (\underset{n\rightarrow \infty}{lim}n[ h(z)-\mathcal{B}_n^{\alpha}(h; z)]\Big)  \right ]\\
&=&\dfrac{1}{2}z(1- z)(fh)''(z)- \dfrac{1}{2}z(1- z){f}''(z)h\left ( z \right )-\dfrac{1}{2}z(1- z){h}''(z)f\left ( z \right )\\
&=&\dfrac{1}{2}z(1- z)\left [ (fh)''(z) -{f}''(z)h\left ( z \right ) -{h}''(z)f\left ( z \right )\right ]=z(1-z){f}'(z){h}'(z).
\end{eqnarray*}
\end{proof}





\subsection{Monotonicity preservation}

\begin{definition}
A system of functions $ \{h_{0},\cdots,h_{n}\}$ is monotonicity preserving % (resp., strictly monotonicity preserving)
 if for any $\nu _{0}\leq \nu _{1} \leq  \cdots\leq \nu_{n} $ %(resp., $\alpha_{0}< \alpha_{1} < \cdots< \alpha_{n} $) 
in $ \mathbb{R} $, the function $ \sum_{i=0}^{n}\nu  _{i}h_{i} $ is increasing. %(resp.,strictly increasing). 
\end{definition}
The Proposition 2.3 of \cite{carnicer1996convexity} outlines the characterization of systems that preserve monotonicity, as presented in the following result. 
\begin{prop}\label{pro1}
Let $\{h_{0},\cdots,h_{n}\}$ be a system of functions defined on an interval $ [a,b]$. Let $ k_{i}:=\sum_{j=i}^{n}h_{j}  $ \ for
$ i \in \left \{  0,1,\cdots,n\right \}$. Then $ \{h_{0},\cdots,h_{n}\}$ is monotonicity preserving if and only if $ k_{0} $ is a constant function and the functions $ k_{i} $ are
increasing for $ i = 1,\cdots,n$. 

\end{prop}
The subsequent result is derived from the preceding proposition.
\begin{prop}\label{pro3.2}
The $\alpha$-Bernstein polynomials $\{F_{n,0}, F_{n,1}, \cdots, F_{n,n}\}$, defined by equations (\ref{formula1}) and (\ref{rec-rel}),  are  monotonicity preserving.
\end{prop}
\begin{proof}
We use the mathematical  induction to verify the result. Consider the $\alpha$-Bernstein polynomials  defined by equation (\ref{rec-rel}) for $n=2$.  According to Proposition \ref{pro1}, the set $\{F_{2,0}, F_{2,1}, F_{2,2}\}$ preserves monotonicity under the conditions that $k_{0}$ is a constant function, and both $k_{1}$ and $k_{2}$ are increasing functions.

Utilizing the partition of unity of $\alpha$-Bernstein polynomials, we obtain
$k_{0}=\sum_{j=0}^{2}F_{2,j}(z)=1$.
On the other hand
\begin{equation*}
\sum_{j=1}^{2}F_{2,j}(z)=z-\alpha(z^{2}-z)\Rightarrow \dfrac{d}{dz} \Big(\sum_{j=1}^{2}F_{2,j}(z)\Big)=1-\alpha(2z-1)\geq 1-\alpha,
\end{equation*}
where $1-\alpha$ takes a non-negative value, therefore $\sum_{j=1}^{2}F_{2,j}(z)$ is an increasing function. Similarly, we show the monotonicity of the last function by taking derivative
\begin{equation*}
\sum_{j=2}^{2}F_{2,j}(z)=z+\alpha(z^{2}-z)\Rightarrow \dfrac{d}{dz} \Big(F_{2,2}(z)\Big)=1+\alpha(2z-1),
\end{equation*}
 where obviously $1+\alpha(2z-1) \geq 0$, so we have the desired result.

Now, we assume that  the $\alpha$-Bernstein polynomials  $\{F_{n-1,0}, F_{n-1,1}, \cdots, F_{n-1,n-1}\}$ are monotonicity preserving, and it remains to conclude the same result for any  $\alpha$-Bernstein polynomials of order $n$,  $\{F_{n,0}, F_{n,1}, \cdots, F_{n,n}\}$.

To do so, we first observe that $k_0=\sum_{j=0}^{n}F_{n,j}(z)=1$, i.e. is a constant function due to  partition of unity. 
For   $k_{i}(z)=\sum_{j=i}^{n}F_{n,j}(z)$, we show the monotonicity by verifying the non-negativity of its  derivative, 
 $ \dfrac{d}{dz}k_{i}(z)=\sum_{j=i}^{n} \dfrac{d}{dz}F_{n,j}(z)$.


Employing relation (\ref{rec-rel}), we have
\begin{eqnarray}
\dfrac{d}{dz}k_{i}(z)&=&-\sum_{j=i}^{n}F_{n-1,j}(z)+(1-z)\sum_{j=i}^{n}\dfrac{d}{dz} F_{n-1,j}(z)+\sum_{j=i}^{n}F_{n-1,j-1}(z)+z\sum_{j=i}^{n}\dfrac{d}{dz}F_{n-1,j-1}(z)\nonumber \\&=&F_{n-1,n}(z)+\left ( -\sum_{j=i}^{n-1}F_{n-1,j}(z)+\sum_{j=i}^{n-1}F_{n-1,j}(z) \right )+F_{n-1,i-1}(z)\nonumber \\&&+(1-z)\sum_{j=i}^{n}\dfrac{d}{dz} F_{n-1,j}(z)+z\sum_{j=i}^{n}\dfrac{d}{dz} F_{n-1,j-1}(z)\nonumber \\&=&F_{n-1,i-1}(z)+(1-z)\sum_{j=i}^{n-1}\dfrac{d}{dz} F_{n-1,j}(z)+z\sum_{j=i-1}^{n-1}\dfrac{d}{dz} F_{n-1,j}(z).\label{11}
\end{eqnarray}
Now, we use the induction hypothesis to see that
\begin{equation}
\sum_{j=i}^{n-1}\dfrac{d}{dz} F_{n-1,j}(z)\geq0, \hspace{1cm}\sum_{j=i-1}^{n-1}\dfrac{d}{dz} F_{n-1,j}(z)\geq0. \label{12}
\end{equation}
and  then it becomes evident that $ \dfrac{d}{dz} k _{i}(z)\geq0$.
\end{proof}


The monotonicity preservation of the $\alpha$-Bernstein operators follows directly from the preceding proposition.


\begin{theorem}
For a function $q$ that is continuous and monotonically increasing on the interval $[0,1]$, the corresponding $\alpha$-Bernstein operators are also increasing.
\end{theorem}
\begin{proof}

Considering the monotonically increasing function $q$, we observe the inequality $q\left(\frac{0}{n}\right) \leq q\left(\frac{1}{n}\right) \leq \cdots \leq q\left(\frac{n}{n}\right)$. Consequently, the result is directly derived from Proposition \ref{pro3.2}.

\end{proof}






\subsection{Convexity preservation}
The preservation of convexity by the $\alpha$-Bernstein operators is verified by the subsequent theorem.

\begin{theorem}
 If $q \in C[0,1] $ is a convex function, then    its $\alpha$-Bernstein operators are all convex.
 
% for $0 \leq \alpha \leq 1$, so are
\end{theorem}

\begin{proof}
We employ the method of induction to establish the result's validity, starting with the base case $n=2$.  Consider a set of real values $ \lambda_{0}, \lambda _{1} , \lambda _{2}$ forming a convex data set, ensuring $\lambda_{2} -2\lambda_{1} +\lambda_{0}\geq 0 $. Our objective is to demonstrate the convexity of the function $ G(z,\alpha)=\sum_{i=0}^{2}\lambda_{i}  F_{2,i}(z)$ over the interval $[0,1]$. Evaluating the second derivative, we get

\begin{equation*}
\dfrac{d^2}{dz^2} G(z,\alpha)=\sum_{i=0}^{2} \lambda_{i}  \dfrac{d^2}{dz^2}  F_{2,i}(z)= 2\alpha \left (\lambda_{2} -2\lambda_{1} +\lambda_{0}   \right )
\end{equation*}

Given this expression and $ \lambda_{2} -2\lambda_{1} +\lambda_{0}\geq 0 $, we can affirm that $ \dfrac{d^2}{dz^2} G(z,\alpha) \geq 0 $.



Presuming the convexity of a set of real values $ \{\beta_{i} \}_{i=0}^{n}$, we take as our induction hypothesis that the expression $ \sum_{i=0}^{n} \beta_{i} F_{n,i}(z) $ is convex.


Now, let's introduce a new set of convex values $ \{\eta_{i}\}_{i=0}^{n+1}$. The goal is to establish the convexity of the function $ \sum_{i=0}^{n+1} \eta_{i} F_{n+1,i}(z) $. To accomplish this, we consider
\begin{equation*} 
G(z,\alpha) =\sum_{i=0}^{n+1} \eta_{i} F_{n+1,i}(z)=\sum_{i=0}^{n+1} \eta_{i}\left [(1-z)F_{n,i}(z)+zF_{n,i-1}(z) \right ],  \end{equation*}
and establish that the function $ \dfrac{d}{dz}G(z,\alpha) $ is  increasing.
\begin{eqnarray*}
\dfrac{d}{dz} G(z,\alpha) &=&\sum_{i=0}^{n+1} \eta_{i} \dfrac{d}{dz} \left [(1-z)F_{n,i}(z)+zF_{n,i-1}(z)  \right ]\\
&=&\sum_{i=0}^{n+1} \eta_{i}\left [-F_{n,i}(z)+(1-z) \dfrac{d}{dz} F_{n,i}(z)+F_{n,i-1}(z)+z \dfrac{d}{dz} F_{n,i-1}(z)  \right ]\\
&=&\sum_{i=0}^{n} ( \eta_{i+1}-\eta_{i})F_{n,i}(z)+(1-z)\sum_{i=0}^{n}\eta_{i} \dfrac{d}{dz} F_{n,i}(z)+z \sum_{i=0}^{n}\eta_{i+1}\dfrac{d}{dz} F_{n,i}(z)
\end{eqnarray*}
Given that $\left \{ \eta_{i}  \right \}_{i=0}^{n+1}$ represents convex data, we observe that $  \eta_{i+1}-2 \eta_{i}+ \eta_{i-1}\geq 0, i=1,\cdots,n-1  $. This inequality implies $  \eta_{i+1}- \eta_{i}\geq \eta_{i}- \eta_{i-1} , i=1,\cdots,n-1  $. In simpler terms, the differences $ \eta_{i}- \eta_{i-1} , i=1,\cdots,n-1 $ form an increasing sequence. By applying Proposition \ref{pro3.2}, we can affirm that the function $ \sum_{i=0}^{n} ( \eta_{i+1}-\eta_{i})F_{n,i}(z) $ is an increasing function.   


Assuming the induction hypothesis and given that $\left \{ \eta_{i} \right \}_{i=0}^{n+1}$ represents convex data, it follows that both $ \sum_{i=0}^{n}\eta_{i}F_{n,i}(z) $ and $ \sum_{i=0}^{n}\eta_{i+1}F_{n,i}(z) $ are convex functions. This implies that their respective derivatives, $ \sum_{i=0}^{n}\eta_{i} \dfrac{d}{dz} F_{n,i}(z)$ and $\sum_{i=0}^{n}\eta_{i+1} \dfrac{d}{dz} F_{n,i}(z) $, are increasing functions.


By the derived results and considering $ z \in [0,1] $, it follows that $ \dfrac{d}{dz} G(z,\alpha) $ is an increasing function. Consequently, this implies $ \dfrac{d^2}{dz^2} G(z,\alpha) \geq 0 $, signifying that $ G(z,\alpha) $ is a convex function.


We observe that if $ q $ is a convex function, then the values $ q\left(\frac{0}{n}\right), q\left(\frac{1}{n}\right), \ldots, q\left(\frac{n}{n}\right) $ form a set of convex data. Consequently, this directly implies that $ R_{n,\alpha}(q;z) $ is also a convex function.
\end{proof}



\section{Conclusion}

The current study makes two significant contributions to the field of approximation theory, particularly concerning Bernstein-like operators.

Firstly, a somewhat generalized formula is introduced for generating Bernstein-like functions. The study illustrates three specific families and opens avenues for constructing new and innovative families.

Secondly, it is demonstrated that the newly presented structure offers a recursive formula for generating $\alpha$-Bernstein polynomials. Utilizing this  representation, several properties are derived, and alternative proofs are provided for established results.


%% The Appendices part is started with the command \appendix;
%% appendix sections are then done as normal sections
%% \appendix

%% \section{}
%% \label{}

%% If you have bibdatabase file and want bibtex to generate the
%% bibitems, please use
%%
%%  \bibliographystyle{elsarticle-num} 
%%  \bibliography{<your bibdatabase>}

%% else use the following coding to input the bibitems directly in the
%% TeX file.


























 \bibliographystyle{elsarticle-num} 
 \bibliography{ref} 





%\begin{thebibliography}{00}

%% \bibitem{label}
%% Text of bibliographic item

%\bibitem{}

%\end{thebibliography}







\end{document}
\endinput
%%
%% End of file `elsarticle-template-num.tex'.
